\chapter{Loss of Smell (Anosmia)}

\section*{Definition}
Complete or partial inability to perceive odors, which can be temporary or permanent, congenital or acquired.

\section*{What Happens in Your Body}
Anosmia results from either conductive loss (nasal obstruction preventing odorants from reaching olfactory epithelium — sinusitis, polyps, URI) or sensorineural loss (damage to olfactory neuroepithelium, olfactory nerve, or central pathways). Viral-induced anosmia (including COVID-19) arises from inflammation and damage to sustentacular cells supporting olfactory neurons.

\section*{Causes}
\begin{itemize}
  \item Viral upper respiratory infection (including COVID-19)
  \item Nasal polyps or chronic rhinosinusitis
  \item Allergic rhinitis
  \item Head trauma with cribriform plate fracture
  \item Aging (presbyosmia)
  \item Smoking or toxic chemical exposure
  \item Neurodegenerative diseases (Parkinson disease, Alzheimer disease)
  \item Nasal tumors or sinus neoplasms
  \item Congenital anosmia (Kallmann syndrome)
  \item Medication side effects (zinc nasal sprays, some antibiotics)
\end{itemize}

\section*{When to See a Doctor}
See your doctor if:
\begin{itemize}
  \item Symptoms are severe or getting worse over time
  \item Sudden anosmia after head injury, anosmia with neurological symptoms, or persistent loss of smell beyond 4 weeks
  \item You have other concerning symptoms such as fever, weight loss, or bleeding
\end{itemize}

\section*{Related Symptoms}
\begin{itemize}
  \item Loss of taste (ageusia)
  \item Nasal congestion
  \item Runny nose
  \item Headache
  \item Neurological symptoms
\end{itemize}
\textbf{Important:} If this symptom persists, your doctor will check for serious underlying causes.

\section*{Self-Care / Home Management}
\begin{itemize}
  \item Perform daily olfactory training using 4 distinct scents (such as rose, lemon, clove, eucalyptus) — sniff each for 10-20 seconds twice daily for at least 12 weeks.
  \item Ensure smoke detectors and natural gas alarms are functional, as anosmia impairs the ability to detect dangerous odours.
  \item Practise good nasal hygiene with saline irrigation to clear nasal passages if congestion is contributing to conductive smell loss.
  \item Add extra flavour (herbs, spices, umami) to food since taste is largely dependent on smell; check food expiration dates carefully to avoid spoiled food.
\end{itemize}

\section*{How Your Doctor Will Evaluate You}
\begin{itemize}
  \item Detailed history including onset (sudden vs gradual), preceding URI or head trauma, nasal symptoms, and neurological symptoms.
  \item Nasal endoscopy to evaluate the olfactory cleft, identify polyps, mucosal inflammation, or structural obstruction.
  \item Objective smell testing using validated tools (University of Pennsylvania Smell Identification Test — UPSIT, Sniffin' Sticks, or Brief Smell Identification Test — B-SIT).
  \item MRI of the brain and olfactory bulbs is indicated when central causes are suspected (tumour, neurodegenerative disease, or following head trauma with anosmia).
\end{itemize}

\section*{Treatment Options}
\begin{itemize}
  \item Olfactory training (daily, structured exposure to 4 scents for 12+ weeks) is the most effective intervention for post-viral and idiopathic anosmia, with 30-60\% recovery rates.
  \item Conductive anosmia from nasal obstruction is treated with intranasal corticosteroids, antihistamines, or surgical polypectomy/septoplasty.
  \item Oral corticosteroids (short course) may be used for severe nasal polyposis or inflammatory sinus disease causing anosmia.
  \item No pharmacological treatment exists for sensorineural anosmia from neurodegenerative disease; management focuses on safety counselling and adaptation strategies.
  \item Medications, lifestyle modifications, and physical therapies may be prescribed as appropriate.
  \item Follow-up ensures treatment effectiveness and allows timely adjustments to the care plan.
\end{itemize}

\section*{References}

\begin{thebibliography}{99}
\bibitem{harrison-anosmia} Longo DL, et al. \emph{Harrison's Principles of Internal Medicine}. 21st ed. McGraw-Hill; 2022. Chapter 35: Olfactory Disorders.
\bibitem{uptodate-anosmia} Doty RL, et al. Evaluation and treatment of anosmia. In: UpToDate, Post TW (Ed), UpToDate, Waltham, MA; 2025.
\bibitem{desir} Desir J, et al. Olfactory dysfunction: a clinical review. \emph{JAMA}. 2023;329(7):565-577.
\bibitem{whitcroft} Whitcroft KL, et al. Olfactory dysfunction: a practical approach. \emph{BMJ}. 2024;385:e078142.
\bibitem{hummel} Hummel T, et al. Olfactory training: clinical practice update. \emph{Rhinology}. 2023;61(2):98-112.
\bibitem{doty-book} Doty RL. \emph{The Neurobiology of Olfaction}. 2nd ed. Cambridge University Press; 2024. Chapter 10: Clinical Olfactory Disorders.
\end{thebibliography}

\clearpage